\pdfbookmark[2]{Zu X.}{nexus-remarks-x}
\begin{center}\large Zu X.\end{center}

Es sei $p$ eine positive Zahl der Form $6n\pm1$, die durch kein Quadrat teilbar ist, und
\[
S_r=\Sigma\left(\frac{s_r}p\right)
\]
wobei $s_r$ alle zu $p$ teilerfremden Zahlen durchlaufen muss, die zwischen $(r-1)\frac p{12}$ und $r\frac p{12}$ liegen. Bezeichnet man die Anzahlen der nicht eigentlich äquivalenten, eigentlich primitiven Formen für die Determinanten $-p$ und $-3p$ mit $C_1$, $C_3$, so findet man leicht (vgl. \textsc{Dirichlet}, Recherches etc., § 11, oder die Note zu VI und VII)
\[
C_1=2(S_1+S_2+S_3),\quad C_3=2(S_1+S_2-S_5-S_6)
\]
oder
\[
C_1=S_1+S_2+S_3+S_4+S_5+S_6,\quad C_3=2(S_2+S_3+S_4+S_5)
\]
je nachdem, ob $p\equiv1$ oder $\equiv3\pmod4$ ist. Berücksichtigt man außerdem, dass
\begin{center}\small
\begin{tabular}{@{}ccccccl@{}}
die&Zahlen&$s_1$ und $s_2$&mit&den&Zahlen&$2s_1$ und $p-2s_6$\\
„&„&$s_3$ und $s_4$&„&„&„&$2s_2$ und $p-2s_5$\\
„&„&$s_5$ und $s_6$&„&„&„&$2s_3$ und $p-2s_4$
\end{tabular}
\end{center}
und ebenso
\begin{center}\small
\begin{tabular}{@{}ccccccl@{}}
die&Zahlen&$s_1,s_2,s_3$&mit&den&Zahlen&$3s_1,\ 3s_5-p,\ p-3s_4$\\
„&„&$s_4,s_5,s_6$&„&„&„&$3s_2,\ 3s_6-p,\ p-3s_3$
\end{tabular}
\end{center}
übereinstimmen und dass im Fall $p\equiv1\pmod4$ die Summe $S_1+S_2+S_3+S_4+S_5+S_6=0$ ist, so erhält man neben den beiden obigen vier weitere Beziehungen zwischen den sechs Summen $S_1,S_2\ldots S_6$, sodass alle aus $C_1$ und $C_3$ bestimmt werden können. Auf diese Weise erhält man, wenn $p\equiv1\pmod4$ ist,
\SourceLeaf{313}
\[
\begin{aligned}
S_1=S_{12}&=\tfrac14\left(1+\left(\frac3p\right)\right)C_1+\tfrac1{12}\left(1+\left(\frac2p\right)\right)C_3\\
S_2=S_{11}&=-\tfrac14\left(1+\left(\frac3p\right)\right)C_1+\tfrac1{12}\left(1+\left(\frac2p\right)\right)C_3\\
S_3=S_{10}&=\tfrac12C_1-\tfrac16\left(1+\left(\frac2p\right)\right)C_3\\
S_4=S_9&=-\tfrac12C_1-\tfrac16\left(2-\left(\frac2p\right)\right)C_3\\
S_5=S_8&=\tfrac14\left(1+\left(\frac3p\right)\right)C_1-\tfrac1{12}\left(5-\left(\frac2p\right)\right)C_3\\
S_6=S_7&=-\tfrac14\left(1+\left(\frac3p\right)\right)C_1+\tfrac1{12}\left(1+\left(\frac2p\right)\right)C_3
\end{aligned}
\]
und, wenn $p\equiv3\pmod4$ ist,
\[
\begin{aligned}
S_1=-S_{12}&=\tfrac1{12}\left(9+3\left(\frac2p\right)-\left(\frac3p\right)+\left(\frac6p\right)\right)C_1-\tfrac14C_3\\
S_2=-S_{11}&=\tfrac14\left(-1+\left(\frac2p\right)+\left(\frac3p\right)-\left(\frac6p\right)\right)C_1+\tfrac14C_3\\
S_3=-S_{10}&=\tfrac16\left(-\left(\frac3p\right)+\left(\frac6p\right)\right)C_1\\
S_4=-S_9&=\tfrac16\left(3-2\left(\frac3p\right)-\left(\frac6p\right)\right)C_1\\
S_5=-S_8&=\tfrac14\left(-1-\left(\frac2p\right)+\left(\frac3p\right)+\left(\frac6p\right)\right)C_1+\tfrac14C_3\\
S_6=-S_7&=\tfrac1{12}\left(3-3\left(\frac2p\right)+\left(\frac3p\right)-\left(\frac6p\right)\right)C_1-\tfrac14C_3
\end{aligned}
\]

Ist $p$ eine Primzahl, so folgen aus dem ersten System die im Text angegebenen Formeln; für die anderen Fälle erhält man ähnliche Formeln aus dem zweiten System.
\begin{flushright}\textsc{R. Dedekind.}\end{flushright}
\BookRule{55mm}
